\section{Introduction}
\label{sec:intro}

The rapid expansion of public blockchains, and of Ethereum in particular,
has reshaped digital finance through decentralised applications and
programmable smart contracts. The same expansion has opened new attack
surfaces. Public reporting attributes more than one billion United States
dollars in losses to cryptocurrency-related scams between 2021 and 2022
alone \cite{chainalysis2023}. These losses arise from a recognisable set
of fraud archetypes: phishing campaigns that exfiltrate private keys
\cite{coinsmart2023}, Ponzi schemes that recycle inbound deposits to
earlier participants \cite{agrawal2023}, smart-contract exploits that
embed malicious behaviour beneath legitimate interfaces, and rug pulls in
which developers withdraw liquidity and abandon their projects
\cite{datavisor2022}.

Defensive techniques from the cryptographic literature, such as threshold
public-key management \cite{mosakheil2023}, multi-signature wallets
\cite{drijvers2019,xiao2022}, and node-authentication schemes
\cite{bao2018,a2chain2020}, reduce the risk of compromise but do not
identify malicious counterparties before a transaction is signed. Earlier
machine-learning detectors trained on tabular features
\cite{ostapowicz2019,blanco2022,ashfaq2022,elmougy2023} achieve respectable
accuracy on snapshots of the data but degrade as adversaries adapt.
Graph-based detectors \cite{kim2023,zhou2021,ctgcn2024} capture
relational anomalies between addresses, yet they overlook the semantic
content of contract bytecode.

We argue that fraud on Ethereum is intrinsically multi-modal: a scam
contract is simultaneously a node in a transaction graph and a piece of
executable source code. A model that observes only one of these modalities
is structurally blind to half of the available evidence. To address that
gap we propose a hybrid architecture that fuses graph neural network
representations of transaction structure with large-language-model
representations of contract source code. Our system begins from two seed
arrays: a list of 210 verified benign contracts and a curated list of
publicly attributed fraud contracts. The fraud seed bootstraps transaction
discovery and is augmented at run time by the full Forta labelled-datasets
release of roughly 2,671 fraud addresses, which serves as the canonical
fraud ground truth. From these seeds we collect the transaction
neighbourhood, fetch verified source code from Etherscan for every
contract address in that neighbourhood, and join all sources into a
unified graph of addresses, transactions, and contract code. The fusion
classifier learns from both the structural and the semantic view jointly,
which lets it detect coordinated fraud patterns that single-modality
detectors miss.

The contributions of this paper are as follows.

\begin{itemize}
    \item We define a needle-first data-collection pipeline that starts
    from two clearly delineated seed arrays of benign and fraud contract
    addresses, expanding outward through the transaction neighbourhood
    rather than scanning the entire chain.
    \item We design a fusion classifier that combines a two-layer
    GraphSAGE encoder with a frozen CodeBERT encoder and report the
    classifier's behaviour on a labelled Ethereum subset of \num{1258}
    verified contracts, \num{183} of which are fraudulent.
    \item We release a multi-file, reproducible implementation of the
    pipeline that replaces a legacy single-notebook prototype, separating
    data collection, preprocessing, modelling, training, evaluation, and
    visualisation into independently testable modules.
    \item We conduct an ablation study comparing the fusion model
    against GNN-only and LLM-only baselines, demonstrating that the
    fusion of modalities outperforms either modality in isolation.
\end{itemize}

The rest of the paper is organised as follows. Section~\ref{sec:motivation}
states our motivation and the limitations of single-modality detectors.
Section~\ref{sec:background} reviews related work along three streams:
graph-centric, code-centric, and hybrid. Section~\ref{sec:methodology}
specifies the architecture and the four-phase data pipeline.
Section~\ref{sec:challenges} discusses the practical challenges encountered
and the countermeasures applied. Section~\ref{sec:prototype} describes
the prototype implementation. Section~\ref{sec:results} reports the
experimental results, including the ablation study. Section~\ref{sec:conclusion}
concludes and outlines future work.
